\documentclass[titlepage]{article}

\title{Free Objects}
\author{Paul Zupan}
\date{Last updated \today}

\begin{document}

\maketitle

\section{Introduction}

In linear algebra, there exists the idea of a basis. Let \(V\) be a vector
space, and suppose that there exists a set of vectors \(\mathcal{B} = \{v_1,
\dots, v_n\}\) in \(V\) such that any vector in \(V\) can be uniquely written
as a linear combination of the vectors in \(\mathcal{B}\). If such a set
exists, we call it a \textit{basis} of \(V\), and we say that \(V\) has
\textit{dimension} \(n\), since \(|\mathcal{B}| = n\).

Free objects generalize bases to other categories, expressing the concept of an
algebraic object being ``freely generated'' by some set.

Let \(\mathcal{C}\) be a category. A \textbf{free object} in \(\mathcal{C}\) is
a triple \((A, X, p)\) where \(A\) is an object.

\end{document}
